\clearpage
\VerdierPhysicalPageAnchor{0184}%
\begingroup\small\setlength{\parskip}{0pt}\setlength{\emergencystretch}{1.5em}
\setlength{\abovedisplayskip}{4pt plus 2pt minus 2pt}%
\setlength{\belowdisplayskip}{4pt plus 2pt minus 2pt}%

\VerdierSemanticAnchor{ch3d-1-3-4-c}
\noindent\textit{c) Tout triangle distingué de $D(\mathcal A)$ est
isomorphe à un triangle distingué du type
$(X,Y,Z,[u],[v],\delta(S))$, où
$S=0\longrightarrow X\xrightarrow{u}Y\xrightarrow{v}Z\longrightarrow0$
est une suite exacte de $\mathcal A$.}

\VerdierSemanticAnchor{ch3d-1-3-4-proof}
La démonstration de a) et b) est laissée au lecteur. Démontrons c). D'après
(chap. II, 2.2.6), tout triangle distingué de $D(\mathcal A)$ est isomorphe
à l'image par $Q\colon K(\mathcal A)\longrightarrow D(\mathcal A)$ d'un
triangle distingué de $K(\mathcal A)$. D'après (chap. I, 3.3.1) et l'axiome
TRII, tout triangle distingué de $K(\mathcal A)$ est isomorphe à un triangle
de la forme :

\VerdierSemanticAnchor{ch3d-1-3-4-cone-triangle}
\[
Y\xrightarrow{\dot{p}(u)}c(\dot u)
\xrightarrow{\dot{q}(u)}X[1]
\xrightarrow{-\dot u[1]}Y[1].
\]
Par suite, tout triangle distingué de $D(\mathcal A)$ est isomorphe à un
triangle du type :

\VerdierSemanticAnchor{ch3d-1-3-4-derived-triangle}
\[
(Y,c(\dot u),X[1],[p(u)],[q(u)],-[u][1]).
\]
Remarquons alors que :

\VerdierSemanticAnchor{ch3d-1-3-4-cone-exact}
\[
S=0\longrightarrow Y\xrightarrow{p(u)}c(u)
\xrightarrow{q(u)}X[1]\longrightarrow0
\]
est une suite exacte de $\operatorname{comp}(\mathcal A)$. Il résulte alors
de (chap. II, 1.2.4) qu'il existe un isomorphisme du triangle :
\[
(Y,c(\dot u),X[1],[p(u)],[q(u)],-[u][1])
\]
sur le triangle :
\[
(Y,c(\dot u),X[1],[p(u)],[q(u)],\delta(S)).
\]

\VerdierSemanticAnchor{ch3d-1-3-5}
\medskip\noindent\textbf{Proposition 1.3.5.}\quad
\textit{a) Tout foncteur
$F\colon\operatorname{comp}(\mathcal A)\longrightarrow C$ qui transforme
les quasi-isomorphismes en isomorphismes de $C$, se factorise de manière
unique par le foncteur
$\operatorname{comp}(\mathcal A)\longrightarrow D(\mathcal A)$.}

\noindent\textit{b) Tout foncteur semi-exact de
$\operatorname{comp}(\mathcal A)$ dans une catégorie abélienne $B$ qui
transforme les quasi-isomorphismes en isomorphismes se factorise d'une
manière unique à travers le $\delta$-foncteur canonique
$\operatorname{comp}(\mathcal A)\longrightarrow D(\mathcal A)$ et un
foncteur cohomologique de $D(\mathcal A)$ dans $B$.}

\VerdierSemanticAnchor{ch3d-1-3-5-proof}
L'assertion b) résulte de l'assertion a) et de la description des triangles
distingués de $D(\mathcal A)$ donnée par (1.3.4). L'assertion a) résulte de
la définition de $D(\mathcal A)$ (chap. II, 2.2.6) et du lemme ci-après.

\VerdierSemanticAnchor{ch3d-1-3-6}
\medskip\noindent\textbf{Lemme 1.3.6.}\quad
\textit{Soient $\mathcal A$ une catégorie additive (non nécessairement
abélienne) et $F$ un foncteur de $\operatorname{comp}(\mathcal A)$ dans
une catégorie $C$. Les conditions suivantes sont équivalentes :}
\endgroup
